\section{Conclusão}

Este trabalho apresentou uma abordagem baseada em modelos de linguagem de grande escala (LLMs) e agentes inteligente para apoiar os desenvolvedores no processo de criação de sistemas, com foco no refinamento de requisitos e na geração de artefatos de desenvolvimento. A proposta buscou explorar como a combinação de multiplos agentes combinado com o conhecimento humano e especializado na area de desenvolvimento de software pode auxiliar na transoformações de ideias iniciais em especificações amis estruturadas e implementaveis de forma mais fácil e rápica.

Esse texto foi escrito utilizando a ferramenta em questão, mostrando que a aplicabilidade da metodologia foi capaz do desenvolvimento completo das principais funcionalidades de forma eficiente. Demonstrando que o estudo e a utilização dessas novas tecnologias vieram para ajudar quando utilizadas da forma correta.

A aplicação experimental demonstrou a sincronia entre humando e máquina, reduzindo a ambiguidades nos requisitos, melhorando a organização das informações e fornecendo suporte durante diferentes etapas o ciclo de desenvolvimento. Vale ressaltar que nesses testes não foram publicados e disponibilizado a ferramenta para uso geral, tendo em vista que não foram utilizados mecanismos de validação de segurança, com isso deve ser utilizado apenas em ambiente local para uso próprio o sistema, ficando de trabalhos futuros a contribuição de como podemos alocar os agentes para atuar como atacantes em busca de vulnerabilidades nela.

Apesar dos resultados incicias indicarem benefícios no uso dessa abordagem, ainda existe desafios relacionados a confiabilidade nas respostas geradas, controle de qualidade e definição de estratégias para validação das implementações realizadas pelos agentes. Dessa forma, o trabalho evidencia a necessidade de novos estudos buscando avaliar esses aspectos em diferentes cenários de desenvolvimento de software.

Com base para trabalhos futuros, pretende-se ampliar a avaliação experimental, analisando métricas relacionadas a qualidade dos requisitos gerados, produtvidade durante o desenvolvimento, interação entre humanos e agentes com novas abordagens, métricas de segurança para avaliar a confiabilidade da aplicação gerada e comparativo da metodologia chaining com a de multiagentes. Além disso, novas integrações com utilizando novos agentes especializados e ferramentas de engenharia de software podem ser investigadas, visando contribuir para o melhor modelo para cada implementação no processo de desenvolvimento.